% GENERATED FILE. Edit canonical lessons, then run npm run generate:books.
% canonical-source-hash: fnv1a64:264812c707268b8a
% canonical-lessons: JA-C463-kitaeru, JA-C463-tairyoku, JA-C463-nerau, JA-C463-hazumu, JA-C463-kisou

\chapter{Training Hard}
\label{ch:ja-training-hard}
\hlchaptermodality{463}

\begin{chapteropening}
\textbf{By the end of this chapter:} \emph{I can talk about training and stamina.}

Complete the last lesson of chapter 463: I can talk about training and stamina.
\end{chapteropening}

\section[\texorpdfstring{\ja{きたえる} (kitaeru)}{きたえる (kitaeru)}]{\ja{きたえる} (kitaeru) — to toughen; to train (the body)}

\label{lesson:JA-C463-kitaeru}



\begin{quote}
Before the new one: say the Japanese for a ranking, then the Japanese for a foul.
\end{quote}

\subsection*{You'll want to know: \ja{きたえる}}
\textbf{\ja{きたえる}} — \emph{kitaeru} — \textquotedblleft{}to toughen; to train (the body)\textquotedblright{}.

\begin{quote}
\textbf{\ja{まいにち} \ja{はしって、からだを} \ja{きたえています。}} — \emph{mainichi hashitte, karada o kitaete imasu} — \textquotedblleft{}I run every day to toughen my body.\textquotedblright{}
\end{quote}

Polite \textbf{\ja{きたえます}}, \emph{kitaemasu}; an ichidan verb. Past (polite): \textbf{\ja{せんしゅが} \ja{からだを} \ja{きたえました。}}, \emph{senshu ga karada o kitaemashita}.

\begin{cousinweb}
\textbf{\ja{きたえる}} (\emph{kitaeru}) is a native Japanese word (wago), inherited rather than borrowed from Chinese; its kanji spelling was attached to it later (JMdict).
\end{cousinweb}

\begin{grammarlens}[title={What you've built}]
One more word for describing what you lived and felt.
\end{grammarlens}

\subsection*{Your turn}
\begin{itemize}
\raggedright
  \item \emph{Say it:} \emph{kitaeru}
  \item \emph{Say it:} \emph{kitaeru}, once more
  \item \emph{Say it:} \emph{mainichi hashitte, karada o kitaete imasu} — I run every day to toughen my body.
  \item \emph{Recall:} say the Japanese for a ranking, then the Japanese for a foul, then say \emph{kitaeru} again
\end{itemize}

\begin{tcolorbox}[breakable,colback=teal!4,colframe=teal!35!black,title={Before you move on}]
What does \ja{きたえる} mean? (\textquotedblleft{}To toughen; to train (the body)\textquotedblright{}.) Say it once more.
\end{tcolorbox}
\section[\texorpdfstring{\ja{たいりょく} (tairyoku)}{たいりょく (tairyoku)}]{\ja{たいりょく} (tairyoku) — stamina; physical strength}

\label{lesson:JA-C463-tairyoku}



\begin{quote}
Before the new one: say the Japanese for a foul, then the Japanese for to toughen.
\end{quote}

\subsection*{You'll want to know: \ja{たいりょく}}
\textbf{\ja{たいりょく}} — \emph{tairyoku} — \textquotedblleft{}stamina; physical strength\textquotedblright{}.

\begin{quote}
\textbf{\ja{かれは} \ja{たいりょくが} \ja{あります。}} — \emph{kare wa tairyoku ga arimasu} — \textquotedblleft{}He has a lot of stamina.\textquotedblright{}
\end{quote}

\begin{cousinweb}
\textbf{\ja{たいりょく}} (\emph{tairyoku}) is Sino-Japanese (kango): two Chinese characters read with their on'yomi, readings borrowed from Chinese. They are \emph{tai}, \textquotedblleft{}body\textquotedblright{} (Middle Chinese \emph{thejX}) and \emph{ryoku}, \textquotedblleft{}power\textquotedblright{} (\emph{lik}) (KANJIDIC2; Middle Chinese in Baxter's transcription).
\end{cousinweb}

\begin{grammarlens}[title={What you've built}]
One more word for describing what you lived and felt.
\end{grammarlens}

\subsection*{Your turn}
\begin{itemize}
\raggedright
  \item \emph{Say it:} \emph{tairyoku}
  \item \emph{Say it:} \emph{tairyoku}, once more
  \item \emph{Say it:} \emph{kare wa tairyoku ga arimasu} — He has a lot of stamina.
  \item \emph{Recall:} say the Japanese for a foul, then the Japanese for to toughen, then say \emph{tairyoku} again
\end{itemize}

\begin{tcolorbox}[breakable,colback=teal!4,colframe=teal!35!black,title={Before you move on}]
What does \ja{たいりょく} mean? (\textquotedblleft{}Stamina; physical strength\textquotedblright{}.) Say it once more.
\end{tcolorbox}
\section[\texorpdfstring{\ja{ねらう} (nerau)}{ねらう (nerau)}]{\ja{ねらう} (nerau) — to aim at}

\label{lesson:JA-C463-nerau}



\begin{quote}
Before the new one: say the Japanese for to toughen, then the Japanese for stamina.
\end{quote}

\subsection*{You'll want to know: \ja{ねらう}}
\textbf{\ja{ねらう}} — \emph{nerau} — \textquotedblleft{}to aim at\textquotedblright{}.

\begin{quote}
\textbf{\ja{せんしゅは} \ja{ゆうしょうを} \ja{ねらっています。}} — \emph{senshu wa yūshō o neratte imasu} — \textquotedblleft{}The players are aiming for the championship.\textquotedblright{}
\end{quote}

Polite \textbf{\ja{ねらいます}}, \emph{neraimasu}; a godan verb. Past (polite): \textbf{\ja{かれが} \ja{まとを} \ja{ねらいました。}}, \emph{kare ga mato o neraimashita}.

\begin{cousinweb}
\textbf{\ja{ねらう}} (\emph{nerau}) is a native Japanese word (wago), inherited rather than borrowed from Chinese; its kanji spelling was attached to it later (JMdict).
\end{cousinweb}

\begin{grammarlens}[title={What you've built}]
One more word for describing what you lived and felt.
\end{grammarlens}

\subsection*{Your turn}
\begin{itemize}
\raggedright
  \item \emph{Say it:} \emph{nerau}
  \item \emph{Say it:} \emph{nerau}, once more
  \item \emph{Say it:} \emph{senshu wa yūshō o neratte imasu} — The players are aiming for the championship.
  \item \emph{Recall:} say the Japanese for to toughen, then the Japanese for stamina, then say \emph{nerau} again
\end{itemize}

\begin{tcolorbox}[breakable,colback=teal!4,colframe=teal!35!black,title={Before you move on}]
What does \ja{ねらう} mean? (\textquotedblleft{}To aim at\textquotedblright{}.) Say it once more.
\end{tcolorbox}
\section[\texorpdfstring{\ja{はずむ} (hazumu)}{はずむ (hazumu)}]{\ja{はずむ} (hazumu) — to bounce}

\label{lesson:JA-C463-hazumu}



\begin{quote}
Before the new one: say the Japanese for stamina, then the Japanese for to aim at.
\end{quote}

\subsection*{You'll want to know: \ja{はずむ}}
\textbf{\ja{はずむ}} — \emph{hazumu} — \textquotedblleft{}to bounce\textquotedblright{}.

\begin{quote}
\textbf{\ja{この} \ja{まりは} \ja{よく} \ja{はずみます。}} — \emph{kono mari wa yoku hazumimasu} — \textquotedblleft{}This ball bounces well.\textquotedblright{}
\end{quote}

Polite \textbf{\ja{はずみます}}, \emph{hazumimasu}; a godan verb. Past (polite): \textbf{\ja{まりが} \ja{たかく} \ja{はずみました。}}, \emph{mari ga takaku hazumimashita}.

\begin{cousinweb}
\textbf{\ja{はずむ}} (\emph{hazumu}) is a native Japanese word (wago), inherited rather than borrowed from Chinese; its kanji spelling was attached to it later (JMdict).
\end{cousinweb}

\begin{grammarlens}[title={What you've built}]
One more word for describing what you lived and felt.
\end{grammarlens}

\subsection*{Your turn}
\begin{itemize}
\raggedright
  \item \emph{Say it:} \emph{hazumu}
  \item \emph{Say it:} \emph{hazumu}, once more
  \item \emph{Say it:} \emph{kono mari wa yoku hazumimasu} — This ball bounces well.
  \item \emph{Recall:} say the Japanese for stamina, then the Japanese for to aim at, then say \emph{hazumu} again
\end{itemize}

\begin{tcolorbox}[breakable,colback=teal!4,colframe=teal!35!black,title={Before you move on}]
What does \ja{はずむ} mean? (\textquotedblleft{}To bounce\textquotedblright{}.) Say it once more.
\end{tcolorbox}
\section[\texorpdfstring{\ja{きそう} (kisou)}{きそう (kisou)}]{\ja{きそう} (kisou) — to compete; to vie}

\label{lesson:JA-C463-kisou}



\begin{quote}
Before the new one: say the Japanese for to aim at, then the Japanese for to bounce.
\end{quote}

\subsection*{You'll want to know: \ja{きそう}}
\textbf{\ja{きそう}} — \emph{kisou} — \textquotedblleft{}to compete; to vie\textquotedblright{}.

\begin{quote}
\textbf{\ja{ふたりの} \ja{せんしゅが} \ja{ゆうしょうを} \ja{きそっています。}} — \emph{futari no senshu ga yūshō o kisotte imasu} — \textquotedblleft{}Two players are competing for the championship.\textquotedblright{}
\end{quote}

Polite \textbf{\ja{きそいます}}, \emph{kisoimasu}; a godan verb. Past (polite): \textbf{\ja{がくせいが} \ja{てんを} \ja{きそいました。}}, \emph{gakusei ga ten o kisoimashita}.

\begin{cousinweb}
\textbf{\ja{きそう}} (\emph{kisou}) is a native Japanese word (wago), inherited rather than borrowed from Chinese; its kanji spelling was attached to it later (JMdict).
\end{cousinweb}

\begin{grammarlens}[title={What you've built}]
One more word for describing what you lived and felt.
\end{grammarlens}

\subsection*{Your turn}
\begin{itemize}
\raggedright
  \item \emph{Say it:} \emph{kisou}
  \item \emph{Say it:} \emph{kisou}, once more
  \item \emph{Say it:} \emph{futari no senshu ga yūshō o kisotte imasu} — Two players are competing for the championship.
  \item \emph{Recall:} say the Japanese for to aim at, then the Japanese for to bounce, then say \emph{kisou} again
\end{itemize}

\begin{tcolorbox}[breakable,colback=teal!4,colframe=teal!35!black,title={Before you move on}]
What does \ja{きそう} mean? (\textquotedblleft{}To compete; to vie\textquotedblright{}.) Say it once more.
\end{tcolorbox}
